% Tutorial set: 06; actual session date: 2026-10-14 (Week 9).
% Source: Week9-tut.pdf pp. 1--3, Q1--Q2 and two linked Colab exercises.
% Colab source inspected on 2026-10-14; no training was executed.
% Q2 values are conditional on the explicitly stated feature conventions.

\begingroup
\renewcommand{\recordid}[1]{}
\exercisegroup{SC4002-TUT06}{Tutorial 06：Linear and Logistic Regression}
\label{tutorial:SC4002-TUT06}

\exercisemeta
  {SC4002-TUT06}
  {2026-10-14}
  {Week9-tut.pdf，pp.~1--3；Q1、Q2 与两个 Colab 编程练习}
  {题干、数值与 Colab 源码已核对；未提供教师答案；未运行训练}
  {已讨论并请求整理；笔记成稿待审阅}

\exercisequestion{SC4002-TUT06-Q01}{Q1：最小二乘拟合电影评分}
\nopagebreak[4]

\textbf{来源与题目（p.~1）。}以影评中 positive words（正面词）数量 $x$ 预测电影评分 $y\in[1,10]$。给定四个样本 $(1,3)$、$(3,6)$、$(5,8)$、$(6,9)$，模型为 $\hat y=wx+b$。用 least squares（最小二乘法）求最优 $w,b$，并预测 $x=4$ 时的评分。

\textbf{对应讲义：}\relatedtopic{SC4002-L08-T01}{线性回归与均方误差}。

\begin{checkedsolution}
最小化平方误差和 $J(w,b)=\sum_{i=1}^{n}(wx_i+b-y_i)^2$。对两个参数求导，令梯度为零：
\[
 \sum_i(wx_i+b-y_i)=0,\qquad
 \sum_i x_i(wx_i+b-y_i)=0.
\]
第一式给出 $b=\bar y-w\bar x$；代回第二式得到
\[
 w=\frac{\sum_i(x_i-\bar x)(y_i-\bar y)}{\sum_i(x_i-\bar x)^2}
 =\frac{n\sum_i x_i y_i-(\sum_i x_i)(\sum_i y_i)}
 {n\sum_i x_i^2-(\sum_i x_i)^2}.
\]
只要 $x_i$ 不全相同，分母就大于零，最优直线唯一。本题
\[
 n=4,\quad \sum x_i=15,\quad \sum y_i=26,\quad
 \sum x_i^2=71,\quad \sum x_i y_i=115,
\]
因此
\[
 \boxed{w=\frac{70}{59}\approx1.1864,\qquad
 b=\frac{121}{59}\approx2.0508,\qquad
 \hat y(4)=\frac{401}{59}\approx6.7966.}
\]
拟合直线经过均值点 $(\bar x,\bar y)=(3.75,6.5)$，可据此自检。四个样本不完全共线，最小二乘拟合无须穿过每个点；将损失除以 $n$ 不改变最优参数。

\begin{center}
% Original plot of the four Q1 samples and their least-squares line.
\begin{tikzpicture}[x=0.86cm,y=0.36cm,>=Latex,font=\small]
  \draw[->] (0,0) -- (7.6,0) node[right] {$x$};
  \draw[->] (0,0) -- (0,11.2) node[above] {$y$};
  \foreach \x in {1,3,4,5,6} {
    \draw (\x,0.15) -- (\x,-0.15) node[below] {\x};
  }
  \foreach \y in {3,6,9} {
    \draw (0.07,\y) -- (-0.07,\y) node[left] {\y};
  }
  \draw[very thick,noteBlue] (0,2.05085) -- (7,10.35593);
  \foreach \x/\y in {1/3,3/6,5/8,6/9} {
    \fill[noteBlue] (\x,\y) circle[radius=2.2pt];
  }
  \draw[dashed,ntuRed] (4,0) -- (4,6.79661);
  \fill[ntuRed] (4,6.79661) circle[radius=2.4pt];
  \node[ntuRed,anchor=north west] at (4.15,5.8) {$\hat y(4)\approx6.80$};
  \node[anchor=west,noteBlue] at (0.8,10) {$\hat y=(70x+121)/59$};
\end{tikzpicture}


{\small 根据 Q1 数据与解析解原创绘制；红点为 $x=4$ 的模型预测。}
\end{center}
\end{checkedsolution}

\exercisequestion{SC4002-TUT06-Q02}{Q2：影评特征、情感概率与交叉熵}
\nopagebreak[4]

\textbf{来源与题目（pp.~1--2）。}给定以下已分词影评：
\begin{quote}
\small
It's hokey . There are virtually no surprises , and the writing is second-rate .
So why was it so enjoyable ? For one thing , the cast is great . Another nice touch
is the music . I was overcome with the urge to get off the couch and start dancing .
It sucked me in , and it'll do the same to you .
\end{quote}
特征依次为：正面词次数、负面词次数、是否含 \texttt{no}、第一／第二人称代词次数、是否含 \texttt{!}、文本长度的对数。完成五项任务：
\begin{enumerate}[label=(\arabic*)]
 \item 提取六维特征向量 $\mathbf x$。
 \item 用 $\mathbf w=[2.5,-5.0,-1.2,0.5,2.0,0.7]^T$、$b=0.1$ 写出二分类逻辑回归函数。
 \item 计算正面概率并给出预测标签。
 \item 若真实情感为正面，计算该样本的损失。
 \item 改为正面、负面、中性三分类，求各类别概率、预测标签及正面真实标签下的损失。给定
\end{enumerate}
\[
 W=\begin{bmatrix}
 1.3&-2.2&-1.0&0.1&0.7&0.5\\
 -2.4&1.9&0.5&-0.4&-1.0&0.2\\
 1.0&0.8&-1.5&1.3&-2.0&0.6
 \end{bmatrix},\qquad
 \mathbf b=\begin{bmatrix}0.1\\0.2\\0.3\end{bmatrix}.
\]

\textbf{对应讲义：}\relatedtopic{SC4002-L08-T02}{逻辑回归与交叉熵}。

\begin{checkedsolution}
\textbf{1. 特征提取口径。}\verify{原题未指定情感词表、长度是否含标点及对数底数。以下数值依赖明确约定，不代表唯一的官方特征答案。}
选用正面词表 \{\texttt{enjoyable}, \texttt{great}, \texttt{nice}\}，负面词表 \{\texttt{hokey}, \texttt{second-rate}\}；按给定空格切分计 66 个 token，包含独立标点，并使用自然对数。得到
\begin{center}
\small
\begin{tabularx}{\textwidth}{@{}c l X@{}}
\toprule
特征 & 数值 & 计数依据 \\
\midrule
$x_1$ & 3 & enjoyable、great、nice 各出现一次。\\
$x_2$ & 2 & hokey、second-rate 各出现一次；本约定不计 sucked。\\
$x_3$ & 1 & 出现 no；此项只记录是否出现。\\
$x_4$ & 3 & I、me、you 各出现一次；包含主格与宾格。\\
$x_5$ & 0 & 没有感叹号。\\
$x_6$ & $\ln66\approx4.189655$ & 若排除独立标点，则为 56 个词，需改用 $\ln56$。\\
\bottomrule
\end{tabularx}
\end{center}
因而 $\mathbf x=[3,2,1,3,0,\ln66]^T$。词表匹配不会自动理解 \texttt{sucked me in} 的语境；更换词表或计数规则后，应重算特征及其下游结果。整体正面的文本仍可包含负面词。

\textbf{2. 二分类得分、概率与标签。}定义正类 $y=1$，则
\begin{align*}
 z&=2.5x_1-5x_2-1.2x_3+0.5x_4+2x_5+0.7x_6+0.1\\
  &=7.5-10-1.2+1.5+0+0.7\ln66+0.1
   \approx0.832758,\\
 p&=P(y=1\mid\mathbf x)=\sigma(z)=\frac1{1+e^{-z}}
   \approx0.696938.
\end{align*}
负面概率为 $1-p\approx0.303062$。采用阈值 $0.5$，预测为\textbf{正面}。

\textbf{3. 二元交叉熵。}单样本损失为
\[
 L=-y\ln p-(1-y)\ln(1-p).
\]
真实标签为正面，即 $y=1$，故 $\boxed{L=-\ln p\approx0.361059}$。分类正确也可以有非零损失；损失取决于分配给真实类别的概率。

\textbf{4. 三分类。}按题目列举顺序，将 $W$ 的三行对应 positive、negative、neutral。计算
\[
 \mathbf z=W\mathbf x+\mathbf b\approx
 \begin{bmatrix}0.994827\\-3.062069\\9.813793\end{bmatrix},\qquad
 p_k=\frac{e^{z_k}}{\sum_{j=1}^{3}e^{z_j}}.
\]
\begin{center}
\begin{tabular}{@{}lrr@{}}
\toprule
类别 & Logit & Softmax 概率 \\
\midrule
positive（正面） & $0.994827$ & $0.000147879$ \\
negative（负面） & $-3.062069$ & $0.000002559$ \\
neutral（中性） & $9.813793$ & $0.999849562$ \\
\bottomrule
\end{tabular}
\end{center}
预测为\textbf{第 3 类 neutral}。真实类别仍为 positive，故
\[
 \boxed{L=-\ln p_{\mathrm{positive}}
 =-\ln\frac{e^{z_1}}{e^{z_1}+e^{z_2}+e^{z_3}}
 \approx8.819116.}
\]
损失较大是因为模型自信地预测了错误类别。损失中应取\emph{真实类别}的概率，不是预测类别的最大概率。两种分类模型使用不同权重，二分类与三分类预测不必一致。
\end{checkedsolution}

\exercisequestion{SC4002-TUT06-C01}{编程练习：用梯度下降拟合直线}
\nopagebreak[4]

\textbf{来源与任务。}原题 p.~2 链接的
\href{https://colab.research.google.com/drive/1KeYjOo36ABjbc-dX1pfJyVyFUQX50dff?usp=sharing}{Linear Regression Colab}。
生成 100 个样本 $x_i\sim\mathcal N(0,1)$、$y_i=5x_i+\epsilon_i$，其中 $\epsilon_i\sim\mathcal N(0,1)$ 为独立随机扰动。已提供线性模型、MSE 损失及学习率 $0.01$ 的 SGD 优化器；补写 1000 个 epoch 的训练循环，每 100 轮打印损失，最后画拟合线。

\textbf{对应讲义：}\relatedtopic{SC4002-L08-T01}{线性回归}、\relatedtopic{SC4002-L08-T03}{梯度下降与批量}。

\begin{checkedsolution}
\texttt{nn.Linear(1, 1)} 实现 $\hat y=wx+b$，只有两个参数；默认 \texttt{nn.MSELoss()} 返回全部 100 个样本的均方误差。参考循环沿用 notebook 中的变量名：
\begingroup\small\linespread{1.05}\selectfont
\begin{verbatim}
for epoch in range(1000):
    optimizer.zero_grad()
    prediction = model(X)
    loss = criterion(prediction, y)
    loss.backward()
    optimizer.step()
    if (epoch + 1) % 100 == 0:
        print(epoch + 1, loss.item())
\end{verbatim}
\endgroup
五个步骤分别是清除旧梯度、前向预测、计算损失、计算梯度和更新参数。\texttt{backward()} 不更新权重；\texttt{step()} 才执行参数更新。每轮使用全部 100 个样本，因此是全批量梯度下降，虽然优化器的类名为 SGD。优化目标与 Q1 相同，这里用迭代法近似求解，而非直接代入最小二乘闭式公式。

\textbf{图与结果的解释。}notebook 提供了散点图和训练前拟合线，并留有训练后的绘图调用。预期拟合线沿散点的主要趋势，$w$ 接近 5、$b$ 接近 0；有限样本的噪声使两者不必精确等于生成参数，损失也不必为零。
\end{checkedsolution}

\exercisequestion{SC4002-TUT06-C02}{编程练习：Iris 三分类与训练步数}
\nopagebreak[4]

\textbf{来源与任务。}原题 p.~3 链接的
\href{https://colab.research.google.com/drive/1C63UPtGeC21-N-TaZftFdWxwMKMLz_Md?usp=sharing}{Logistic Regression Colab}。
使用 Iris 数据，输入四个测量值（花萼长／宽、花瓣长／宽），判断三个鸢尾花品种。根据已给配置补写训练循环，记录损失随 iteration 的变化，最后计算测试准确率。

\textbf{对应讲义：}\relatedtopic{SC4002-L08-T02}{Softmax 与交叉熵}、\relatedtopic{SC4002-L08-T03}{Epoch 与批量更新}。

\begin{checkedsolution}
\textbf{1. 数据与模型。}源码读入 150 个样本，删除 Id 列，将 Species 编码为 setosa$=0$、versicolor$=1$、virginica$=2$。以 \texttt{test\_size=0.25}、\texttt{random\_state=42} 划分为 112 个训练样本和 38 个测试样本。特征转换为浮点张量，标签转换为整数类别索引。

模型为 \texttt{nn.Linear(4, 3)}，即 $\mathbf z=W\mathbf x+\mathbf b$，其中 $W\in\mathbb R^{3\times4}$、$\mathbf b\in\mathbb R^3$，共 15 个参数。一个大小为 $B$ 的批次将 $B\times4$ 的输入映射为 $B\times3$ 的 logits，每行对应一个样本。类别编号没有大小或距离含义。

\textbf{2. 损失接口。}\texttt{nn.CrossEntropyLoss()} 接收原始 logits 和整数标签，内部完成等价于 LogSoftmax 加负对数似然的计算。因此训练时不提前加 Softmax：
\[
 L_{\mathrm{batch}}=-\frac1B\sum_{i=1}^{B}
 \ln\frac{e^{z_{i,y_i}}}{\sum_{k=0}^{2}e^{z_{i,k}}}.
\]
接口依据 \href{https://docs.pytorch.org/docs/stable/generated/torch.nn.CrossEntropyLoss.html}{PyTorch CrossEntropyLoss 文档}；源码使用 SGD，学习率为 $0.001$。

\textbf{3. 原代码的训练次数不一致。}notebook 先用 \texttt{batch\_size=100} 和 \texttt{n\_iters=10000} 算出
\[
 \texttt{num\_epochs}=\left\lfloor\frac{10000}{112/100}\right\rfloor=8928,
\]
之后创建 DataLoader 时却改为 \texttt{batch\_size=16}。实际每轮有 $112/16=7$ 个 batch；若直接训练 8928 轮，会更新 $8928\times7=62496$ 次，而非 10000 次。Epoch 表示遍历一次训练集，iteration 表示一次参数更新。为匹配变量 \texttt{n\_iters} 的意图，可直接按更新次数终止：
\begingroup\small\linespread{1.05}\selectfont
\begin{verbatim}
iteration_list, loss_list = [], []
step = 0
model.train()
while step < n_iters:
    for x_batch, y_batch in train_loader:
        optimizer.zero_grad()
        logits = model(x_batch)
        loss = error(logits, y_batch)
        loss.backward()
        optimizer.step()
        step += 1
        iteration_list.append(step)
        loss_list.append(loss.item())
        if step >= n_iters:
            break
\end{verbatim}
\endgroup
该参考循环记录各次更新前的 mini-batch loss；不同 batch 的难度不同，曲线可有波动，不要求每一步都下降。

\textbf{4. 测试与解释。}测试时停止参数更新，取每行得分的最大类别：
\begingroup\small\linespread{1.05}\selectfont
\begin{verbatim}
model.eval()
with torch.no_grad():
    logits = model(testX)
    predicted = logits.argmax(dim=1)
    accuracy = (predicted == testY).float().mean().item()
\end{verbatim}
\endgroup
Softmax 保持大小次序，因此 $\arg\max_k z_k=\arg\max_k\operatorname{softmax}(\mathbf z)_k$。只需预测类别时可以直接取 logits 的最大值；要展示概率时再计算 Softmax。准确率为测试样本中预测正确的比例，分母为 38。

训练损失反映拟合过程，测试准确率用于检查泛化，不应据测试结果反复选择训练轮数。源码只固定了数据划分的随机状态，未固定 PyTorch 初始化与批次打乱的随机状态，重新运行的训练结果可能变化。
\end{checkedsolution}
\endgroup
